\Needspace{8\baselineskip}
\section{分块的数据复用与三种策略}
\label{l8:sec:tiling}
\subsection{访问计数、算术强度与复用机会}
对于不接触边界的一个一维输出，直接实现读取 $K$ 个输入与 $K$ 个系数，共 $2K$ 次标量读取，并执行 $K$ 次乘法和 $K$ 次加法，共 $2K$ 次浮点运算。二维 $K\times K$ 滤波器对应 $2K^2$ 次标量读取与 $2K^2$ 次浮点运算。因此增加滤波器宽度时，直接算法的读取量和计算量一起增长，其复用前的算术强度较低。

\term{算术强度}{arithmetic intensity}用计算量与数据搬运量之比描述计算对带宽的需求。按每个 \mbox{\code{float}} 为 $4$ B，将标量读取次数换成字节：一维内部输出读入 $8K$ B、写出 $4$ B。忽略复用时，仅以读入量为分母得到 $2K/(8K)=0.25$ FLOP/B；把输出写入也计入，得到
\begin{equation}
 I_{\mathrm{direct,1D}}=\frac{2K}{8K+4},\qquad
 I_{\mathrm{direct,2D}}=\frac{2K^2}{8K^2+4}.
\label{l8:eq:ai}
\end{equation}
较低的算术强度说明直接算法容易受到带宽限制；常量缓存与共享内存的作用，就是减少同一份数据在较慢存储层次上的重复获取。

\textbf{计数口径。}这里的“读取次数”首先指代码层面的标量访问。缓存命中、常量同地址访问以及成组内存事务会改变实际 DRAM 流量，所以不能直接把 $2K$ 次读取等同于 $2K$ 次独立 DRAM 事务。边界线程又会在直接核函数中跳过无效项，此时实际执行的读取与乘加比内部输出少。后面的精确例题始终先说明统计的是输入请求、字节还是乘加次数。

\subsection{输出分块、输入覆盖范围与 halo}
设一个线程块生成 $B$ 个连续输出，块号为 $b$，输出起点为 $s=bB$。输出范围是 $[s,s+B-1]$。最左输出需要从 $s-r$ 开始读取，最右输出需要读到 $s+B-1+r$，因此整个块需要的输入范围为
\begin{equation}
 [s-r,\ s+B+r-1],\qquad
 T_{\mathrm{in}}=B+2r=B+K-1.
\label{l8:eq:tilewidth}
\end{equation}
这里 $B$ 是输出块宽度，$T_{\mathrm{in}}$ 是包含边缘邻域的输入块宽度。输入核心区与输出对齐，左右各多出的 $r$ 个位置称为\term{晕区}{halo}。

\begin{figure}[H]
\centering
\includegraphics[width=.87\textwidth]{tiles_halos}
\caption{$W=16$、$B=4$、$K=5$ 时的四个输出块及其输入邻域。中间块的 halo 是真实输入，首末块还包含越界的幽灵位置。}
\label{l8:fig:tiles}
\end{figure}

以 $b=1$ 为例，这个块生成 $P[4]$ 到 $P[7]$，输入需要 $N[2]$ 到 $N[9]$。核心区是 $N[4:7]$，左 halo 是 $N[2:3]$，右 halo 是 $N[8:9]$。相邻块的\textbf{输出不重叠、输入可以重叠}，所以线程块可以独立写各自的输出，但边缘输入可能被多个块读取。第一个块需要的 $N[-2]$ 与 $N[-1]$ 才属于全局数组外的幽灵位置，按边界规则填零。

\textbf{块内复用的数量。}对于 $B=4,K=5$，四个输出分别需要五个输入，共产生 $20$ 次输入使用，但覆盖的不同输入只有八个。$N[2]$ 到 $N[9]$ 的块内使用次数依次为 $1,2,3,4,4,3,2,1$，总和为 $20$。一旦把这八个值放入共享内存，块内后续使用就可以从共享内存读取。

对于远离全局边界的完整块，直接方法的输入请求数为 $BK$；把完整输入块载入一次则为 $B+2r$。因此只比较 $N$ 的块内读取请求，可写出复用因子
\begin{equation}
 R_N=\frac{BK}{B+2r}=\frac{BK}{B+K-1}.
\label{l8:eq:reuse}
\end{equation}
这给出示例中的 $20/8=2.5$。固定 $K$ 而增大 $B$ 时，halo 相对占比 $2r/B$ 下降，$R_N$ 趋近于 $K$。这些是输入覆盖关系推得的请求量比较，没有计入系数缓存、同步或硬件资源对实际速度的影响。

\subsection{线程数量与 halo 放置的组合}
分块之后仍有两个决定：线程数与输出块匹配，还是与输入块匹配；halo 是否也放入共享内存。由此形成三种主要策略。

\begin{table}[H]
\centering\small
\begin{tabularx}{\textwidth}{@{}p{.17\textwidth}Y Y Y@{}}
\toprule
比较项 & 策略 1 & 策略 2 & 策略 3 \\
\midrule
线程块覆盖 & $B$ 个输出 & $B+2r$ 个输入 & $B$ 个输出 \\
每块线程数 & $B$ & $B+2r$ & $B$ \\
共享内存内容 & 核心区与两侧 halo & 核心区与两侧 halo & 仅核心区 \\
载入方式 & 部分线程额外载入 & 每线程载入一个输入 & 每线程载入一个核心输入 \\
计算方式 & 每线程计算一个输出 & 仅部分线程计算输出 & 每线程计算一个输出 \\
计算时 halo & 从共享内存读取 & 从共享内存读取 & 从全局地址读取 \\
主要取舍 & 载入阶段任务不均 & 计算阶段部分线程闲置 & 计算循环选择两种读取路径 \\
\bottomrule
\end{tabularx}
\caption{以相同的 $B$ 个输出为比较对象，线程数与输入存放方式的差异。}
\end{table}

\begin{figure}[H]
\centering
\includegraphics[width=.98\textwidth]{three_strategies}
\caption{三种分块策略的输入、输出和线程工作范围。策略 1 分步载入完整输入块；策略 2 为输入配置线程；策略 3 将 halo 留给全局读取路径。}
\end{figure}

\textbf{策略 1 的取舍。}将 halo 一并载入共享内存后，卷积计算循环的输入都来自同一种位置，不需要在每次乘加时选择共享或全局路径。连续位置可以按组载入，也有利于合并访问。代价是共享内存由 $B$ 个元素增加到 $B+2r$ 个元素，而且部分线程在载入阶段必须多做工作。载入任务按线程编号分配时，可能出现部分线程执行分支、其他线程不执行的情况。

\textbf{策略 3 的取舍。}只缓存核心区，把共享内存优先用于块内复用较高的输入。计算每一个乘加项时判断输入是否位于本块核心区：位于核心区就读共享内存，位于 halo 就按全局地址读取。同一 warp 内的线程可能选择不同路径；窄 halo 的访问也可能只利用一次成组内存传输的一小部分。

按 $r\le B$ 的完整内部块展开计数，左侧各输出对左 halo 的读取数为 $r,r-1,\ldots,1$，右侧对称，因此两侧一共 $r(r+1)$ 次全局 halo 请求。再加 $B$ 次核心区载入，策略 3 对 $N$ 的请求数为 $B+r(r+1)$。在 $B=4,r=2$ 的例子中，这个数为 $10$，策略 1 则是 $8$。该比较按程序请求计数；策略 3 的 halo 请求可能由缓存满足，不能只靠这两个数判定实测速度。

\textbf{策略 2 的取舍。}让每个线程载入一个输入，载入阶段无需由少量线程额外承担 halo；随后仅有与输出对应的线程执行乘加。按相同输出宽度 $B$ 比较，计算输出的线程比例为 $B/(B+2r)$。这种分配减少高延迟载入阶段的工作分歧，同时让部分线程不参与较低延迟的计算阶段。全局数组首末处的边界填零仍然需要处理。
\FloatBarrier
